\section*{الفصل \ref{ch:b1:rate-laws} --- الحركية الكيميائية: قوانين السرعة}

\begin{solution}{exo:b1:rate-laws:1}
$v = \frac14 \times 2.0 \times 10^{-4} = \qty{5.0e-5}{mol.L^{-1}.s^{-1}}$؛
ويتكوّن ثنائي الأكسجين بالسرعة نفسها، \qty{5.0e-5}{}؛ ويختفي \ce{N2O5} بالسرعة
$2v = \qty{1.0e-4}{mol.L^{-1}.s^{-1}}$.
\end{solution}

\begin{solution}{exo:b1:rate-laws:2}
الرتبة 0: \unit{mol.L^{-1}.s^{-1}}؛ الرتبة 1: \unit{s^{-1}}؛ الرتبة 2:
\unit{L.mol^{-1}.s^{-1}}؛ الرتبة 3: \unit{L^2.mol^{-2}.s^{-1}}.
\end{solution}

\begin{solution}{exo:b1:rate-laws:3}
$t_{1/2} = \ln2/k = 0.693/(3.0 \times 10^{-3}) = \qty{231}{s}$. وبعد
\qty{600}{s}، $\eu^{-kt} = \eu^{-1.8} = 0.165$: يبقى 16.5\,\%.
\end{solution}

\begin{solution}{exo:b1:rate-laws:4}
\omterm{def:b1:rate-laws:half-life}{زمن نصف تفاعل} متناسب مع $1/a_0$ يعني الرتبة 2. وإذا كان تنصيف $a_0$
ينصّف \omterm{def:b1:rate-laws:half-life}{زمن نصف التفاعل}، فإن $t_{1/2} \propto a_0$: الرتبة 0.
\end{solution}

\begin{solution}{exo:b1:rate-laws:5}
$[\ce{A}]$ ليست خطية في $t$ (الخسائر 0.0259، 0.0192، 0.0142، 0.0106).
$\ln[\ce{A}] = -2.303$، $-2.602$، $-2.902$، $-3.202$، $-3.503$: ينخفض
بمقدار 0.300 كل \qty{10}{min}، فهو مستقيم. الرتبة 1، $k =
\qty{0.0300}{min^{-1}}$.
\end{solution}

\begin{solution}{exo:b1:rate-laws:6}
مضاعفة $[\ce{A}]_0$ تضرب $v_0$ في $4 = 2^2$: الرتبة 2 بالنسبة إلى \ce{A}.
ومضاعفة $[\ce{B}]_0$ تضربها في 2: الرتبة 1 بالنسبة إلى \ce{B}. $v =
k[\ce{A}]^2[\ce{B}]$ مع $k = 2.0 \times 10^{-6}/(0.010^2 \times 0.010) =
\qty{2.0}{L^2.mol^{-2}.s^{-1}}$.
\end{solution}

\begin{solution}{exo:b1:rate-laws:7}
\ce{B} بإفراط مئة ضعف: فيبقى تركيزه
\qty{0.50}{mol/L}، و$v = k'[\ce{A}]$ مع $k' = k[\ce{B}]_0$. ومنه
$k = 0.012/0.50 = \qty{0.024}{L.mol^{-1}.s^{-1}}$.
\end{solution}

\begin{solution}{exo:b1:rate-laws:8}
$E_a = R\ln5/(1/300 - 1/320) = 8.314 \times 1.609/(2.083 \times 10^{-4}) =
\qty{64.2}{kJ/mol}$. وعند \qty{310}{K}: $\ln k = \ln(1.0 \times 10^{-3}) +
\frac{E_a}{R}\big(\frac{1}{300} - \frac{1}{310}\big) = -6.908 + 0.831$،
$k = \qty{2.3e-3}{s^{-1}}$.
\end{solution}

\begin{solution}{exo:b1:rate-laws:9}
هنا $a = 2$: $1/[\ce{A}] = 1/a_0 + 2kt$، $t_{1/2} = 1/(2k a_0) = 1/(2
\times 0.50 \times 0.020) = \qty{50}{s}$. وعند \qty{100}{s}: $1/[\ce{A}] = 50
+ 100 = \qty{150}{L/mol}$، $[\ce{A}] = \qty{6.7e-3}{mol/L}$.
\end{solution}

\begin{solution}{exo:b1:rate-laws:10}
$[\ce{A}] = a_0(1 - f) = a_0\eu^{-akt_f}$ يعطي $t_f = -\ln(1-f)/(ak)$،
وبأزمنة نصف التفاعل $t_f/t_{1/2} = -\ln(1-f)/\ln2$: 1 من أجل 50\,\%، و2 من أجل 75\,\%،
و3.32 من أجل 90\,\%، و9.97 من أجل 99.9\,\%: نحو عشرة أزمنة نصف تفاعل كي يكون التفاعل
«منتهيًا».
\end{solution}

\begin{solution}{exo:b1:rate-laws:11}
$m(t) = 10 - 0.40\,t$ (\unit{mg}، \unit{h}): الرتبة 0. \omterm{def:b1:rate-laws:half-life}{زمن نصف التفاعل}
$10/(2 \times 0.40) = \qty{12.5}{h}$؛ وتفرغ بعد \qty{25}{h}. فالجرعة
المحرَّرة في الساعة لا تتعلق بما تبقى: إمداد منتظم، وهذا
هو الغرض من اللصيقة.
\end{solution}

\begin{solution}{exo:b1:rate-laws:12}
$E_a = R\ln2/(1/298.15 - 1/308.15) = 8.314 \times 0.693/(1.088 \times
10^{-4}) = \qty{52.9}{kJ/mol}$. وبين 273.15 و\qty{283.15}{K}:
$\exp\big(\frac{52900}{8.314}(\frac{1}{273.15} - \frac{1}{283.15})\big) =
2.3$: \omterm{def:b1:rate-laws:activation-energy}{فطاقة التنشيط} نفسها تعطي عاملًا أكبر عند درجة حرارة
أدنى.
\end{solution}

\begin{solution}{pb:b1:rate-laws:1}
\textbf{1.} الخسائر 8.6 و7.9 و7.2 و6.5 نقطة: ليست ثابتة، فليست
الرتبة 0.
\textbf{2.} $\ln$: 4.605، 4.515، 4.425، 4.335، 4.245؛ ينخفض بمقدار 0.090
كل 10 أيام: مستقيم، الرتبة 1.
\textbf{3.} $k_{60} = 0.090/10 = \qty{9.0e-3}{d^{-1}}$.
\textbf{4.} $t_{1/2} = \ln2/k_{60} = \qty{77}{d}$.
\textbf{5.} $100\,\eu^{-0.54} = 58.3\,\%$.
\textbf{6.} $t_{90} = \ln(10/9)/k_{60} = 0.1054/0.0090 = \qty{11.7}{d}$.
\textbf{7.} يبقى 90\,\%: $\eu^{-kt_{90}} = 0.9$، فيكون $t_{90} = \ln(10/9)/k$.
\textbf{8.} $t_{90}/t_{1/2} = \ln(10/9)/\ln2 = 0.152$.
\textbf{9.} في الرتبة 1 لا يتعلق الكسر المتبقي إلا بالمقدار $kt$، لا
بالكمية الابتدائية.
\textbf{10.} عند \qty{25}{\celsius} يستغرق التفكك سنوات؛ والتسخين
يسرّعه (أرهينيوس) فيمكن قياسه في أسابيع.
\textbf{11.} في الرتبة 2 يكون \omterm{def:b1:rate-laws:half-life}{زمن نصف التفاعل} $1/(ak\,a_0)$: فالقرص المضاعف
القوة يتفكك أسرع بمرتين نسبيًا ويدوم مدة أقل.
\textbf{12.} $1/T$ (\unit{K^{-1}}): \num{3.1934e-3}، \num{3.0945e-3}،
\num{3.0017e-3}؛ $\ln k$: $-6.669$، $-5.661$، $-4.711$.
\textbf{13.} الميل $= (-4.711 + 6.669)/(3.0017 \times 10^{-3} - 3.1934 \times
10^{-3}) = \qty{-1.021e4}{K}$.
\textbf{14.} $E_a = 8.314 \times 1.021 \times 10^4 = \qty{85}{kJ/mol}$.
\textbf{15.} القيمة المتوقعة $\ln k_{50} = -4.711 - 1.021 \times 10^4 \times (3.0945 -
3.0017) \times 10^{-3} = -5.658$، والمقيسة $-5.661$: النقطة على المستقيم.
\textbf{16.} $A = k_{60}\,\eu^{E_a/RT} = 0.0090\,\eu^{10215/333.15} =
\qty{1.9e11}{d^{-1}}$.
\textbf{17.} $\exp\big(1.021 \times 10^4(\frac{1}{298.15} -
\frac{1}{308.15})\big) = 3.0$: أكثر من العامل 2 في القاعدة، لأن
$E_a$ أكبر من \qty{53}{kJ/mol}.
\textbf{18.} $k_{30} = k_{60}\exp\big(-1.021 \times 10^4(\frac{1}{303.15} -
\frac{1}{333.15})\big) = \qty{4.3e-4}{d^{-1}}$.
\textbf{19.} $t_{90} = 0.1054/4.33 \times 10^{-4} = \qty{243}{d}$، أي 8.0
أشهر.
\textbf{20.} $1 - \eu^{-3 \times 0.0090} = 2.7\,\%$.
\textbf{21.} تتناقص مدة الصلاحية بسرعة مع درجة الحرارة (8 أشهر عند
\qty{30}{\celsius} مقابل 14 عند \qty{25}{\celsius}): فتاريخ انتهاء الصلاحية
لا يصح إلا دون \qty{25}{\celsius}.
\textbf{22.} $k_{25} = k_{60}\exp\big(-1.021 \times 10^4(\frac{1}{298.15} -
\frac{1}{333.15})\big) = \qty{2.46e-4}{d^{-1}}$.
\textbf{23.} $t_{90} = 0.1054/2.46 \times 10^{-4} = \qty{428}{d}$، أي
\textbf{14 شهرًا}: وهو التاريخ المطبوع على العلبة، المستخرج من ستة
أسابيع من القياسات.
\end{solution}
